
\documentclass{article}
\usepackage{colortbl}
\usepackage{makecell}
\usepackage{multirow}
\usepackage{supertabular}

\begin{document}

\newcounter{utterance}

\centering \large Interaction Transcript for game `cladder', experiment `full\_v1.5\_default', episode 771 with llama{-}3.1{-}8b{-}instruct.
\vspace{24pt}

{ \footnotesize  \setcounter{utterance}{1}
\setlength{\tabcolsep}{0pt}
\begin{supertabular}{c@{$\;$}|p{.15\linewidth}@{}p{.15\linewidth}p{.15\linewidth}p{.15\linewidth}p{.15\linewidth}p{.15\linewidth}}
    \# & \multicolumn{2}{c}{Player} && \multicolumn{2}{c}{Game Master} \\
    \hline

    \theutterance \stepcounter{utterance}  
    & & & \multicolumn{4}{p{0.6\linewidth}}{
        \cellcolor[rgb]{0.9,0.9,0.9}{
            \makecell[{{p{\linewidth}}}]{
                \texttt{\tiny{[P1$\langle$GM]}}
                \texttt{You are an expert in causal inference. The following question is not a typical commonsense query, but rather a meticulously designed question created by a professor specializing in causal inference, intended to assess the students' mastery of the course content.} \\
\texttt{Imagine a self{-}contained, hypothetical world with only the following conditions, and without any unmentioned factors or causal relationships: Gender has a direct effect on occupation and salary. Skill has a direct effect on occupation and salary. Occupation has a direct effect on salary. Skill is unobserved. For individuals who are not male and blue{-}collar workers, the probability of high salary is 26\%. For individuals who are not male and white{-}collar workers, the probability of high salary is 61\%. For individuals who are male and blue{-}collar workers, the probability of high salary is 55\%. For individuals who are male and white{-}collar workers, the probability of high salary is 68\%. For individuals who are not male and with low skill levels, the probability of white{-}collar job is 40\%. For individuals who are not male and with high skill levels, the probability of white{-}collar job is 15\%. For individuals who are male and with low skill levels, the probability of white{-}collar job is 2\%. For individuals who are male and with high skill levels, the probability of white{-}collar job is 24\%. The overall probability of high skill level is 11\%. If we disregard the mediation effect through occupation, would gender negatively affect salary?} \\
\texttt{Start your answer with 'Yes' or 'No', followed by additional reasoning or evidence to support your explanation.} \\
            }
        }
    }
     \\ \\

    \theutterance \stepcounter{utterance}  
    & \multicolumn{4}{p{0.6\linewidth}}{
        \cellcolor[rgb]{0.561,0.737,0.561}{
            \makecell[{{p{\linewidth}}}]{
                \texttt{\tiny{[P1$\rangle$GM]}}
                \texttt{No.} \\
\\ 
\texttt{To assess the effect of gender on salary while disregarding the mediation effect through occupation, we can use the backdoor criterion. This involves identifying a set of covariates that block all backdoor paths from gender to salary.} \\
\\ 
\texttt{First, let's identify the backdoor paths from gender to salary. The direct effect of gender on salary is given. The backdoor paths from gender to salary are:} \\
\\ 
\texttt{1. Gender → Skill → Salary} \\
\texttt{2. Gender → Occupation → Salary} \\
\\ 
\texttt{We are given the following conditional probabilities:} \\
\\ 
\texttt{{-} P(Salary|Not Male, Not Blue{-}collar) = 0.26} \\
\texttt{{-} P(Salary|Not Male, White{-}collar) = 0.61} \\
\texttt{{-} P(Salary|Male, Blue{-}collar) = 0.55} \\
\texttt{{-} P(Salary|Male, White{-}collar) = 0.68} \\
\\ 
\texttt{We can use these probabilities to estimate the effect of gender on salary while disregarding the mediation effect through occupation.} \\
\\ 
\texttt{Let's consider the following:} \\
\\ 
\texttt{{-} P(White{-}collar|Not Male, Low Skill) = 0.40} \\
\texttt{{-} P(White{-}collar|Not Male, High Skill) = 0.15} \\
\texttt{{-} P(White{-}collar|Male, Low Skill) = 0.02} \\
\texttt{{-} P(White{-}collar|Male, High Skill) = 0.24} \\
\\ 
\texttt{We are also given that P(High Skill) = 0.11.} \\
\\ 
\texttt{Using the law of total probability, we can estimate the probability of high salary for individuals with high skill levels:} \\
\\ 
\texttt{P(Salary|High Skill) = P(Salary|High Skill, Male)P(Male|High Skill) + P(Salary|High Skill, Not Male)P(Not Male|High Skill)} \\
\\ 
\texttt{We know that P(Salary|High Skill, Male) = 0.68 and P(Salary|High Skill, Not Male) = 0.61. We also know that P(High Skill) = 0.11, so P(Male|High Skill) = 0.24 and P(Not Male|High Skill) = 0.76.} \\
\\ 
\texttt{Therefore, P(Salary|High Skill) = (0.68)(0.24) + (0.61)(0.76) = 0.65.} \\
\\ 
\texttt{Now, let's consider the probability of high salary for individuals with low skill levels:} \\
\\ 
\texttt{P(Salary|Low Skill) = P(Salary|Low Skill, Male)P(Male|Low Skill) + P(Salary|Low Skill, Not Male)P(Not Male|Low Skill)} \\
\\ 
\texttt{We know that P(Salary|Low Skill, Male) = 0.55 and P(Salary|Low Skill, Not Male) = 0.26. We also know that P(Low Skill) = 0.89, so P(Male|Low Skill) = 0.46 and P(Not Male|Low Skill) = 0.54.} \\
            }
        }
    }
    & & \\ \\

    \theutterance \stepcounter{utterance}  
    & & & \multicolumn{2}{p{0.3\linewidth}}{
        \cellcolor[rgb]{0.9,0.9,0.9}{
            \makecell[{{p{\linewidth}}}]{
                \texttt{\tiny{[GM$|$GM]}}
                \texttt{No} \\
            }
        }
    }
    & & \\ \\

    \theutterance \stepcounter{utterance}  
    & & & \multicolumn{2}{p{0.3\linewidth}}{
        \cellcolor[rgb]{0.9,0.9,0.9}{
            \makecell[{{p{\linewidth}}}]{
                \texttt{\tiny{[GM$|$GM]}}
                \texttt{no} \\
            }
        }
    }
    & & \\ \\

    \theutterance \stepcounter{utterance}  
    & & & \multicolumn{2}{p{0.3\linewidth}}{
        \cellcolor[rgb]{0.9,0.9,0.9}{
            \makecell[{{p{\linewidth}}}]{
                \texttt{\tiny{[GM$|$GM]}}
                \texttt{game\_result = WIN} \\
            }
        }
    }
    & & \\ \\

\end{supertabular}
}

\end{document}
